% !TEX program = lualatex
% =============================================================================
%  ЗАПИТ НА ПУБЛІЧНУ ІНФОРМАЦІЮ — Закон України «Про доступ до публічної
%  інформації», ст. 19 ч. 5: ім'я запитувача, поштова або електронна адреса
%  (і телефон, якщо є); опис інформації чи реквізити документа; підпис і дата.
%  Строк відповіді — 5 робочих днів (ст. 20). Форма довільна (ст. 19 ч. 4).
% =============================================================================
\documentclass[14pt]{dstu4163}

\dstusetup{
  blank            = none,
  doctype          = {Запит},
  doctype-subtitle = {на отримання публічної інформації},
}

\addressee{%
  Умовна міська рада

  вул. Зразкова, 1, м. Умовне, 00000}

\applicant{%
  Запитувач:\\
  Андрій \surname{Мельничук}\\
  вул. Січових Стрільців, 10, кв. 5,\\
  м. Київ, 04053\\
  тел. +380 00 000 00 00\\
  e-mail: a.melnychuk@example.com}

\begin{document}
\makeletterhead

Відповідно до статей 19 і 20 Закону України «Про доступ до публічної
інформації» прошу надати:

1.~Знеособлені дані про кількість договорів оренди комунального
майна за 2023–2025 роки.

2.~Пояснення щодо методики розрахунку орендної плати
за комунальне майно.

Інформацію прошу надати в електронній формі у форматі CSV або XLSX
на адресу електронної пошти a.melnychuk@example.com.

\signature[date = 18.11.2026]{}{Андрій \surname{Мельничук}}

\end{document}
